\chapter{Resumen}

En esta tesina se aborda un problema práctico vinculado al proceso de escrutinio provisorio con asistencia digital de la Provincia de Santa Fe: detectar automáticamente errores de confección en telegramas electorales verificando que las firmas manuscritas de las autoridades de mesa estén presentes y sean consistentes a lo largo de todas las páginas del documento. A diferencia del problema clásico de \textit{Verificación Automática de Firmas Manuscritas} (VAFM), aquí no se busca detectar falsificaciones sino inconsistencias de autoría entre firmas que deberían corresponder a la misma persona, con el objetivo de reducir retransmisiones y trabajo humano innecesario en el centro de cómputo \cite{unr-informe,tesis-parodi}. 

El problema se formula como una tarea de clasificación binaria donde, dados dos recortes de firmas, se debe decidir si fueron realizadas por el mismo autor o no \cite{metricas-binarias}. Para ello se estudian conceptos fundamentales de Aprendizaje Automático y Aprendizaje Profundo, prestando especial atención al diseño correcto de los conjuntos de entrenamiento, validación y test, así como a la elección de métricas robustas al desbalance y acordes al uso práctico del sistema (\textit{precision}, \textit{recall}, curvas ROC y Precision–Recall, etc.). La tesina contrasta enfoques basados en extracción manual de características y clasificadores clásicos (como el esquema SVM + grilla circular de \cite{tesis-parodi}) con métodos modernos de redes neuronales convolucionales y arquitecturas siamesas para aprendizaje de representaciones.

En la fase experimental se parte de la confección de un conjunto de datos propio a partir de telegramas reales: se recortan automáticamente las zonas de firmas aprovechando que la ubicación es fija, y se entrena primero un clasificador para distinguir entre “hay firma” y “no hay firma”, que sirve como herramienta para limpiar y estructurar el \textit{dataset} definitivo de pares de firmas. Luego se entrenan distintos modelos de similitud: un \textit{baseline} inspirado en la tesis de Parodi (features manuales + Random Forest/SVM) y diversos modelos de redes siamesas convolucionales, variando arquitectura, función de costo y estrategia de generación de pares y tripletas (incluyendo variantes \textit{online} y \textit{offline} de \textit{triplet loss}). 

Los resultados muestran que las arquitecturas siamesas convolucionales, entrenadas con funciones de costo orientadas a aprendizaje de métricas (como \textit{triplet loss}), superan de manera clara a los métodos basados en extracción manual de características tanto en conjuntos de datos estándar de firmas como en el dominio específico de telegramas electorales. Se alcanza una región de operación con \textit{recall} cercano a 0.9 y alta \textit{precision}, consistente con el objetivo práctico de reducir a una fracción las retransmisiones debidas a errores de firmas, manteniendo a la vez un número manejable de falsas alarmas. Finalmente se discuten las limitaciones del enfoque (dependencia del conjunto de datos, supuestos sobre el formato de los telegramas, ausencia de falsificaciones intencionales) y se señalan líneas de trabajo futuro, entre ellas el uso de modelos generativos para enriquecer las representaciones y la aplicación del método a otros problemas de similitud de imágenes.


